(1) 设界面切向为 \(x\)，Snell 条件要求 \(k_x=k_0n_1\sin\theta\) 在两侧相同。透射介质满足 \(k_x^2+k_{z,2}^2=k_0^2n_2^2\)，因此在 \(\theta>\theta_c\) 时
\[
k_{z,2}=ik_0\kappa,\qquad
\kappa=\sqrt{n_1^2\sin^2\theta-n_2^2}>0.
\]
透射场振幅随 \(z>0\) 按 \(\exp(-k_0\kappa z)\) 衰减，其 \(1/e\) 振幅穿透深度为
\[
\delta_{\rm ev}=\frac1{k_0\kappa}=\frac{\lambda}{2\pi\kappa}.
\]
不应把这一深度与强度的 \(1/e\) 深度（后者为一半）混淆。

(2) TE 振幅反射系数可写为 \(r_s=(k_{z,1}-k_{z,2})/(k_{z,1}+k_{z,2})\)，其中 \(k_{z,1}=k_0n_1\cos\theta\)。代入虚数法向波数：
\[
r_s=\frac{n_1\cos\theta-i\kappa}{n_1\cos\theta+i\kappa}
=e^{i\varphi_s},
\quad \varphi_s=-2\arctan\frac{\kappa}{n_1\cos\theta}.
\]
分母是分子的复共轭，因而 \(|r_s|=1\)，表示无吸收时反射的法向平均能流与入射相同，但透射介质仍可存在非零倏逝场。

(3) 对窄角谱光束，沿 \(x\) 的反射相位 \(\varphi_s(k_x)\) 导致波包包络驻相位置产生位移
\[
\Delta_{\rm GH}=-\frac{d\varphi_s}{dk_x}
=-\frac{1}{k_0n_1\cos\theta}\frac{d\varphi_s}{d\theta}.
\]
对上式直接求导并使用 \(\kappa^2=n_1^2\sin^2\theta-n_2^2\)，得到 \(d\varphi_s/d\theta=-2n_1\sin\theta/\kappa\)，因此
\[
\Delta_{\rm GH}=\frac{2\tan\theta}{k_0\kappa}>0 .
\]
在这一 TE、平面界面和窄角谱约定中，位移沿光束的界面投影前进方向。

(4) 当 \(\theta\downarrow\theta_c\)，\(\kappa\to0\)，穿透深度和理想驻相位移都趋大；这并非预言有限宽度光束真的产生无限位移，因为角谱可能跨越全反射临界角，窄角谱近似失效。全反射时透射侧均匀倏逝波在 \(z\) 方向的周期平均坡印廷通量为零，但沿界面方向可有非零局域能量流，二者不矛盾。
